\documentclass[11pt,a4paper]{article}
\usepackage[T1]{fontenc}
\usepackage[utf8]{inputenc}
\usepackage{amsmath,amssymb,amsthm}
\usepackage[margin=1in]{geometry}
\usepackage[colorlinks=true,linkcolor=blue,urlcolor=blue]{hyperref}
\theoremstyle{plain}
\newtheorem{theorem}{Theorem}
\newtheorem{lemma}[theorem]{Lemma}
\newtheorem{proposition}[theorem]{Proposition}
\newtheorem{corollary}[theorem]{Corollary}
\theoremstyle{definition}
\newtheorem{remark}[theorem]{Remark}

\title{Recovering the unwritten recurrences of the table OEIS A237481}
\author{Adrian Perez Fontelles\\ \small Independent researcher}
\date{2 September 2026}
\begin{document}
\maketitle

\begin{abstract}
OEIS A237481 is a two-dimensional table $T(n,k)$ whose empirical recurrences are, for several of
its lines, given only as an ORDER --- ``k=2: [order 30]'' and the like --- with the
coefficients in a linked file rather than in the entry. Those conjectures can still be settled,
and the recurrences recovered. A line of the table is a fixed-width or fixed-height array
count, hence a walk count on finitely many vertices, hence satisfies some monic recurrence of
order at most that number; Berlekamp--Massey on twice as many exact terms returns the MINIMAL
one. Where its order equals the order the entry states --- and where the same holds after the
first few terms are dropped --- any recurrence of that order the line satisfies has a
characteristic polynomial that is a multiple of the minimal one and of equal degree, hence is
that polynomial. This note settles column 2, column 3.
\end{abstract}

\noindent\small 2020 Mathematics Subject Classification. 05A15, 11B37, 15A18.\normalsize

\section{The table and the unwritten conjectures}

OEIS A237481 is ``T(n,k)=Number of (n+1)X(k+1) 0..2 arrays with the maximum plus the lower median of every 2X2 subblock equal.''. It is read by antidiagonals and begins
\[
81, 331, 331, 1337, 2529, 1337, 5819, 19175,\ \dots
\]
The entry carries, among its empirical claims:
\begin{quote}\small
k=2: [order 30]\\
k=3: [order 73]
\end{quote}
As of the ``Last modified'' line on the live entry (Jul 23 2025 10:12:34, revision 6) these are still
recorded as empirical, and the recurrences themselves are not in the entry text.

\section{Why a line of the table is a walk count}

Fix the index that the line fixes. The entry defines $T(n,k)$ as a number of arrays of a shape
determined by $n$ and $k$, subject to a condition local to a bounded window of consecutive
lines. Substituting the fixed index into the entry's own wording turns the definition into an
ordinary one-parameter array count, and for such a count the admissible windows are the
vertices of a finite digraph and an array is a walk. So the line is a walk count on $S$
vertices, and by the Cayley--Hamilton theorem it satisfies a monic integer recurrence of order
at most $S$.

\section{Recovering the recurrences}

\begin{lemma}\label{lem:bm}
A sequence known to satisfy some linear recurrence of order at most $S$ is determined, with its
minimal recurrence, by its first $2S$ terms; Berlekamp--Massey applied to them returns that
minimal recurrence exactly.
\end{lemma}

\subsection*{Column $k=2$}
Substituting the fixed index into the entry's wording gives an ordinary one-parameter array
count, a walk on $S=69$ vertices. Berlekamp--Massey on $2S=138$ exact terms
returns a minimal recurrence of order $30$ --- the order the entry states --- with
integer coefficients
\[
(c_1,\dots,c_{30})=(17,\ -4,\ -1270,\ 4614,\ 32689,\ -189159,\ -315240,\ 3345462,\ -441205,\ -31268255,\ 33370256,\ 163213626,\ -289548851,\ -461018671,\ 1232093408,\ 554566403,\ -2950288660,\ 316531223,\ 4092259150,\ -1771232392,\ -3233425862,\ 2126933890,\ 1373888416,\ -1177913072,\ -282612928,\ 315961632,\ 21654656,\ -36843008,\ -182272,\ 1425408).
\]
so that $a(m)=\sum_{i=1}^{30}c_i\,a(m-i)$ for every $m$. Removing the first
$30+4$ terms and repeating gives order $30$ again, which is what the
identification below needs.

\subsection*{Column $k=3$}
Substituting the fixed index into the entry's wording gives an ordinary one-parameter array
count, a walk on $S=178$ vertices. Berlekamp--Massey on $2S=356$ exact terms
returns a minimal recurrence of order $73$ --- the order the entry states --- with
integer coefficients
\[
(c_1,\dots,c_{73})=(24,\ 216,\ -8276,\ -1861,\ 1166736,\ -3191279,\ -91235323,\ 402333178,\ 4480606213,\ -25766386496,\ -147151963954,\ 1058614532664,\ 3317970084353,\ -30543061732853,\ -50598457068202,\ 650242307521613,\ 461230172021974,\ -10540791350189129,\ -558909097815716,\ 132937549899031809,\ -57223965741744522,\ -1324468730941933534,\ 1068475609904223217,\ 10539801712410825456,\ -11762454331311616768,\ -67512288405742533507,\ 93552613049106259400,\ 349841627389919200693,\ -572731685896476840451,\ -1470149915261604564869,\ 2778362090494465163199,\ 5008537685104250658538,\ -10849881033010974200915,\ -13780131245888880250708,\ 34423586908151572264928,\ 30341252295210587029509,\ -89201129692217400669204,\ -52472198117878266766311,\ 189265217629044248029334,\ 68429978480049972687950,\ -328946443664790007682551,\ -60142235348902158951529,\ 467597635368710546414862,\ 18536338333016940955471,\ -541795401357281641764800,\ 41609493021211604296382,\ 508925504515936762736291,\ -88067228082807638562910,\ -384524615318262807103153,\ 97809824091929438892961,\ 231151296852062265828100,\ -75236200373596837632797,\ -108878797974707354158960,\ 42614253121620477412880,\ 39314843400498432782476,\ -18032354966914111409296,\ -10525600761768982851384,\ 5669544682403681996848,\ 1973425202866123192608,\ -1298509864416787685856,\ -228864027886370438912,\ 209007471775592362240,\ 9771309398512790272,\ -22252005104655600640,\ 1208604555681060864,\ 1403067397584431104,\ -197351936450617344,\ -40186121476718592,\ 10066000183050240,\ -39609195233280,\ -151273089269760,\ 13741521371136,\ -363746820096).
\]
so that $a(m)=\sum_{i=1}^{73}c_i\,a(m-i)$ for every $m$. Removing the first
$73+4$ terms and repeating gives order $73$ again, which is what the
identification below needs.

\begin{theorem}
For each line above, the displayed recurrence holds for every index, and it is the recurrence
the entry records.
\end{theorem}

\begin{proof}
It holds by Lemma~\ref{lem:bm}. For the identification, the entry asserts that the line
satisfies SOME recurrence of the stated order, possibly only from some index on. Let $q$ be its
characteristic polynomial and $q_{\min}$ the minimal polynomial of the tail. Then
$q_{\min}\mid q$, and the second Berlekamp--Massey run --- on the line with its first few
terms removed --- shows $\deg q_{\min}$ equals the stated order, which is $\deg q$. Both are
monic, so $q=q_{\min}$.
\end{proof}

\section{Verification}

Three checks.

First, each line's model was compared against the entry's own published values for that line,
taken from the antidiagonal DATA read in either orientation and from the ``Table starts''
block, and agrees at every available term. This is the only point at which reading the entry's
English enters, and a misreading gives different counts.

Second, each recovered recurrence was evaluated directly on the model's values, with no
Berlekamp--Massey involved, and holds at every index tested.

Third, each order was computed twice by independent routes --- modulo a $61$-bit prime, where
every intermediate is a machine word, and again in exact rational arithmetic --- and once more
on the line with its first few terms dropped, which is what the identification requires. All
agree.

\begin{thebibliography}{9}
\bibitem{oeis} The OEIS Foundation, \emph{The On-Line Encyclopedia of Integer
Sequences}, \url{https://oeis.org}, 2026. Sequence A237481.
\bibitem{massey} J.~L.~Massey, Shift-register synthesis and BCH decoding, \emph{IEEE
Trans. Inform. Theory} \textbf{15} (1969), 122--127.
\bibitem{stanley} R.~P.~Stanley, \emph{Enumerative Combinatorics}, Volume 1, 2nd ed.,
Cambridge University Press, 2011. (Transfer-matrix method, Section 4.7.)
\bibitem{horn} R.~A.~Horn and C.~R.~Johnson, \emph{Matrix Analysis}, 2nd ed., Cambridge
University Press, 2013.
\end{thebibliography}

\end{document}
